% L4b student companion deck.
%
% The notebook is the live teaching document; the slides are the students'
% concise note-taking companion. Match its topic and example sequence. The disclaimer is intentionally
% slide 2, by instructor preference.
\documentclass[aspectratio=169,11pt]{beamer}
\usepackage{vnslides}

\newcommand{\Pmeasure}{\mathbb{P}}
\newcommand{\E}{\mathbb{E}}
\newcommand{\Var}{\operatorname{Var}}
\newcommand{\Cov}{\operatorname{Cov}}
\newcommand{\mug}{\mu_g}
\newcommand{\Dacc}{\mathcal{D}_{T,0}(g_y)}
\newcommand{\Ddis}{\mathcal{D}_{T,0}^{-1}(g_y)}

\title{\texorpdfstring{{\vnheavy L4b}}{L4b}: Single Asset Geometric Brownian Motion Models}
\subtitle{CHEME 5660 Cornell University}
\author{Copyright Jeffrey Varner 2026. All Rights Reserved}

\begin{document}

\vntitlepage

\begin{frame}{Disclaimer and Risks}
  \vspace{0.6em}
  \textbf{\ink{Training only.}} This content is for training and informational
  purposes only. The teaching team does not make, give, or endorse any offer or
  solicitation to buy or sell securities or derivative products, or provide
  investment or trading advice or strategy.

  \vspace{1.1em}
  \textbf{\ink{Risk.}} Trading involves risk and is generally not recommended for
  those with limited resources, experience, or a low risk tolerance. Past
  performance does not guarantee future success.

  \vspace{1.1em}
  \textbf{\ink{Responsibility.}} You are responsible for your investment decisions
  based on your financial circumstances, objectives, risk tolerance, and
  liquidity needs.
\end{frame}

\begin{frame}{Objectives for Today}
  By the end of this lecture, you can:

  \vspace{0.8em}
  \begin{itemize}\setlength{\itemsep}{0.55em}
    \item \hilite{Describe and simulate GBM}: Explain drift and random
      fluctuations, interpret the price distribution, and simulate prices
      using the exact solution.
    \item \hilite{Estimate and interpret GBM parameters}: Estimate mean growth
      and volatility from historical prices, recover the drift, and estimate
      mean growth using log-price regression.
    \item \hilite{Assess GBM and calculate trade probabilities}: Compare GBM
      with observed growth-rate behavior and calculate the probability of
      exceeding a target scaled NPV at a scheduled sale time.
  \end{itemize}
\end{frame}

\begin{frame}{Lectures and Examples}
  \textbf{\ink{Lecture:}}
  \href{https://github.com/varnerlab/CHEME-5660-CourseRepository-Fall-2026/blob/main/lectures/week-4/L4b/CHEME-5660-L4b-Lecture-SingleAsset-GeometricBrownianMotion-TradeRule-Fall-2026.ipynb}{Single-Asset Geometric Brownian Motion}

  \vspace{1.1em}
  \textbf{\ink{Example:}}
  \href{https://github.com/varnerlab/CHEME-5660-CourseRepository-Fall-2026/blob/main/lectures/week-4/L4b/CHEME-5660-L4b-Example-Parameters-SAGBM-Fall-2026.ipynb}{Estimating GBM Parameters}

  Estimate mean growth rate and volatility from historical prices, recover the
  arithmetic drift, and compare simulated price paths with the data.

  \vspace{1.1em}
  \textbf{\ink{Example:}}
  \href{https://github.com/varnerlab/CHEME-5660-CourseRepository-Fall-2026/blob/main/lectures/week-4/L4b/CHEME-5660-L4b-Example-GBM-NPV-TradeRule-Fall-2026.ipynb}{GBM Terminal-Target Probabilities}

  Compute the probability that a long stock position exceeds a target scaled
  NPV at a scheduled sale time.
\end{frame}

\begin{frame}{Company Profile: Jane Street}
  \href{https://www.janestreet.com/}{Jane Street} is a privately held
  \hilite{quantitative trading firm}. Like Citadel Securities in L4a, it
  uses mathematical models and software in trading.

  \vspace{0.4em}
  \begin{itemize}\setlength{\itemsep}{0.5em}
    \item \hilite{Trading business}: Jane Street trades mainly with its own
      capital and provides liquidity to institutional clients. Citadel Securities
      also executes orders for retail brokers and institutional investors.
    \item \hilite{Products}: Jane Street is known for market making in
      exchange-traded funds (ETFs), which let investors trade baskets of assets.
      Both firms operate across multiple asset classes.
    \item \hilite{Technology}: Jane Street uses OCaml, a functional programming
      language, in its trading systems.
  \end{itemize}

  \vspace{0.4em}
  A trade's outcome depends on price changes before the sale. We will
  model this uncertainty with GBM and calculate the probability of exceeding
  a target scaled NPV.
\end{frame}

\begin{frame}{From Lattices to Continuous Prices}
  In L4a, lattice models used discrete time steps and a finite set of
  price movements. What changes as we reduce the time step?

  \vspace{0.4em}
  \begin{itemize}\setlength{\itemsep}{0.5em}
    \item Keep the holding period fixed and divide it into more steps.
    \item Shrink the up and down price movements and adjust their probabilities
      so the mean and variance of the holding-period log return remain
      approximately unchanged.
    \item For suitably scaled, independent binomial steps, the accumulated
      log return $\ln(S_T/S_0)$ approaches a normal distribution.
  \end{itemize}

  \vspace{0.5em}
  Since $S_T=S_0\exp[\ln(S_T/S_0)]$, the terminal price approaches a
  lognormal distribution. The corresponding continuous-time price model is
  \hilite{geometric Brownian motion}.

  \vspace{0.4em}
  We introduce this model through its stochastic differential equation,
  then use its solution to calculate a trade's target probability.
\end{frame}

\begin{frame}{Single Asset Geometric Brownian Motion (GBM)}
  \href{https://gwern.net/doc/economics/1965-samuelson.pdf}{Paul Samuelson (1965)}
  used geometric Brownian motion to model stock prices. The model became
  widely known through
  \href{https://www.cs.princeton.edu/courses/archive/fall08/cos323/papers/black_scholes73.pdf}{Black--Scholes--Merton (BSM) option pricing}.
  Random fluctuations scale with the current price, and prices remain positive.

  \vspace{0.5em}
  Let $S(t)>0$ be the share price (USD/share), with $t$ and the holding period
  $T$ measured in years. The constant parameters are drift $\mu\in\mathbb R$
  (year$^{-1}$) and volatility $\sigma>0$ (year$^{-1/2}$).
  With Wiener increment $dW(t)$, GBM is given by:
  \[
    \frac{dS(t)}{S(t)}
    =\underbrace{\mu\,dt}_{\text{drift}}
    +\underbrace{\sigma\,dW(t)}_{\text{noise}},
    \qquad 0\le t\le T.
  \]

  \vspace{0.5em}
  The drift sets the expected proportional change per unit time. The Wiener
  increment supplies the random shock. Let's define this noise process.
\end{frame}

\begin{frame}{The Wiener Process}
  A Wiener process $\{W(t),\,0\le t\le T\}$ is a continuous real-valued
  stochastic process with three properties:

  \vspace{0.6em}
  \begin{itemize}\setlength{\itemsep}{0.6em}
    \item The process starts at zero: $W(0)=0$ with probability one.
    \item Increments over disjoint intervals
      $W(t_1)-W(t_0),\ldots,W(t_k)-W(t_{k-1})$ are independent for any
      $0\le t_0<t_1<\cdots<t_k\le T$.
    \item Each increment is normal with mean zero and variance equal to the
      elapsed time, $W(t)-W(s)\sim\mathcal{N}(0,t-s)$ for $0\le s<t\le T$.
  \end{itemize}

  \vspace{0.5em}
  The third property lets us write any increment as a scaled standard normal:
  \[
    W(t)-W(s)=\sqrt{t-s}\;Z,\qquad Z\sim\mathcal{N}(0,1).
  \]
\end{frame}

\begin{frame}{The Analytical Solution}
  Take $t_0=0$ and let $S_0$ be today's price. Solving the SDE with It\^o
  calculus gives the share price at time $T$:
  \[
    \boxed{S_T=S_0\exp\!\Bigl[
      \underbrace{\bigl(\mu-\tfrac{\sigma^2}{2}\bigr)T}_{\text{mean log growth}}
      +\underbrace{\sigma\sqrt{T}\,Z}_{\text{shock}}\Bigr],
    \qquad Z\sim\mathcal{N}(0,1).}
  \]

  Thus the log price ratio $\ln(S_T/S_0)$ is normal with mean
  $(\mu-\sigma^2/2)T$ and variance $\sigma^2T$, and the price $S_T$ is
  \hilite{lognormal} and strictly positive.

  \vspace{0.4em}
  It\^o calculus gives the \hilite{half-variance correction} $-\sigma^2/2$.
  This makes the mean growth rate differ from the drift $\mu$. Next, we
  compute the mean and variance of future prices.

  \slidenote{Derivation:}{\href{https://github.com/varnerlab/CHEME-5660-CourseRepository-Fall-2026/blob/main/lectures/week-4/L4b/CHEME-5660-L4b-GBM-Solution-Derivation-Fall-2026.ipynb}{Deriving the GBM Solution}}
\end{frame}

\begin{frame}{Mean and Variance of the GBM Price}
  From the lognormal moments, the expectation and variance of $S_T$ are:
  \[
    \underbrace{\E(S_T)}_{\text{mean price}}=S_0e^{\mu T},
    \qquad
    \underbrace{\Var(S_T)}_{\text{price variance}}
    =S_0^2e^{2\mu T}\bigl(e^{\sigma^2T}-1\bigr).
  \]

  Properties:
  \begin{itemize}\setlength{\itemsep}{0.4em}
    \item The mean averages over possible GBM outcomes. It grows exponentially
      when $\mu>0$ and declines when $\mu<0$, but individual price paths can
      rise or fall around this expectation.
    \item The variance describes the spread of prices:
      $S_0^2e^{2\mu T}$ is the squared expected price, while $e^{\sigma^2T}-1$
      captures the effect of volatility over the holding period.
  \end{itemize}

  \vspace{0.5em}
  For fixed initial price, drift, and holding period, increasing
  volatility widens the price distribution without changing the expected price.
\end{frame}

\begin{frame}{Discrete-Time GBM Model}
  To simulate a price path, use a grid $t_j=j\Delta t$, $j=0,1,\ldots,N$,
  with time step $\Delta t$ (years) and horizon $T=N\Delta t$. Applying the
  GBM solution between successive times gives the \hilite{one-step transition}:
  \[
    \boxed{S_{t_j}=S_{t_{j-1}}\exp\!\Bigl[
      \bigl(\mu-\tfrac{\sigma^2}{2}\bigr)\Delta t
      +\sigma\sqrt{\Delta t}\,Z_j\Bigr].}
  \]

  Here, $j=1,\ldots,N$, and $Z_1,\ldots,Z_N$ are \hilite{independent}
  standard normal random variables. Their independence follows from the
  independent Wiener increments over successive intervals.

  \vspace{0.6em}
  For constant $\mu$ and $\sigma$, this transition gives exact GBM prices at
  the grid points. We use it to generate price paths with Monte Carlo simulation.
\end{frame}

\begin{frame}{Monte Carlo Simulation of GBM}
  Given $S_0$, $\mu$, $\sigma$, $\Delta t$, $N$ steps, and $M$ paths, store the
  simulated prices in $\mathbf{S}\in\mathbb{R}^{(N+1)\times M}$, one path per
  column, rows indexed by $j=0,\ldots,N$.

  \vspace{0.5em}
  For each sample path (each column of $\mathbf{S}$):
  \begin{enumerate}\setlength{\itemsep}{0.45em}
    \item Set the price in row $j=0$ to $S_0$.
    \item For $j=1,\ldots,N$: draw $Z_j\sim\mathcal{N}(0,1)$ and apply the
      one-step transition:
      \[
        S_{t_j}\gets S_{t_{j-1}}\exp\!\Bigl[
          \bigl(\mu-\tfrac{\sigma^2}{2}\bigr)\Delta t
          +\sigma\sqrt{\Delta t}\,Z_j\Bigr].
      \]
  \end{enumerate}

  \vspace{0.3em}
  Each step of each path uses a fresh independent draw.
  The result is $M$ price trajectories of $N$ steps each.

  \vspace{0.4em}
  Let's estimate the drift and volatility needed for simulation from
  historical share prices.
\end{frame}

\begin{frame}{Estimation of GBM Parameters}
  To connect GBM to data, divide the one-step transition by
  $S_{t_{j-1}}$, take logarithms, and divide by $\Delta t$. This gives the
  \hilite{continuously compounded growth rate} $g_j$ (year$^{-1}$) used in L3a:
  \[
    \boxed{g_j\equiv\frac{1}{\Delta t}\ln\frac{S_{t_j}}{S_{t_{j-1}}}
    =\underbrace{\bigl(\mu-\tfrac{\sigma^2}{2}\bigr)}_{\text{mean growth }\mug}
    +\underbrace{\frac{\sigma}{\sqrt{\Delta t}}}_{\text{growth-rate std}}\,Z_j.}
  \]

  Because the $Z_j$ are independent standard normals, growth rates over equally
  spaced, non-overlapping steps are independent and normal with mean
  $\mug\equiv\mu-\sigma^2/2$ and variance $\sigma^2/\Delta t$.

  \vspace{0.3em}
  \begin{itemize}\setlength{\itemsep}{0.4em}
    \item \hilite{Drift versus growth}: $\mu$ is the SDE drift; the mean growth
      rate $\mug$ is estimated from a growth-rate series. They differ by the
      half-variance correction.
    \item \hilite{Volatility from dispersion}: the growth-rate standard deviation
      is $\sigma/\sqrt{\Delta t}$, while the one-step log-return standard
      deviation is $\sigma\sqrt{\Delta t}$. We use either to estimate $\sigma$.
  \end{itemize}
\end{frame}

\begin{frame}{Volatility}
  From $N+1$ prices we form $N$ growth rates $g_1,\ldots,g_N$. Their
  sample mean $g'$ and sample standard deviation $\sigma_g$ (our L3a risk
  measure) are:
  \[
    g'=\frac1N\sum_{j=1}^Ng_j,
    \qquad
    \underbrace{\sigma_g}_{\text{growth-rate std}}
    =\sqrt{\frac{1}{N-1}\sum_{j=1}^N\bigl(g_j-g'\bigr)^2}.
  \]

  Estimating the mean leaves $N-1$ degrees of freedom. To recover GBM
  volatility from $\sigma_g$, use $\Var(g_j)=\sigma^2/\Delta t$:
  $\sigma_g$ estimates $\sigma/\sqrt{\Delta t}$, giving the volatility estimate:
  \[
    \boxed{\hat\sigma=\sigma_g\sqrt{\Delta t}=\sigma_{\mathrm{ann}}.}
  \]

  This matches L3a's annualized volatility $\sigma_{\mathrm{ann}}$, with units
  year$^{-1/2}$.

  For daily data $\Delta t=1/252$: divide the daily growth-rate standard
  deviation by $\sqrt{252}$, or equivalently multiply the daily log-return
  standard deviation by $\sqrt{252}$.
\end{frame}

\begin{frame}{Mean Growth and Drift}
  With volatility estimated, we need mean growth to recover the drift.
  Since $\mu=\mug+\sigma^2/2$, an estimate $\hat{\mu}_g$ gives:
  \[
    \underbrace{\hat\mu}_{\text{SDE drift}}
    =\underbrace{\hat{\mu}_g}_{\text{mean growth}}
    +\underbrace{\tfrac12\hat\sigma^2}_{\text{correction}}.
  \]

  Two common ways to estimate $\mug$:
  \begin{itemize}\setlength{\itemsep}{0.55em}
    \item \hilite{Method 1, average growth rate}: $\hat{\mu}_g=g'$, the
      sample mean and maximum likelihood estimate under GBM.
    \item \hilite{Method 2, linear regression}: fit log price as a line in time.
      The slope estimates $\mug$.
  \end{itemize}

  \vspace{0.4em}
  Let's estimate mean growth by regressing log price on time.
\end{frame}

\begin{frame}{Linear Regression: The Model}
  Suppose we observe prices $S_{t_0},\ldots,S_{t_N}$ at times
  $t_j=j\Delta t$. Taking the logarithm of the GBM solution gives:
  \[
    \ln S_{t_j}
    =\underbrace{\ln S_0}_{\text{intercept}}
    +\underbrace{\mug\,t_j}_{\text{trend}}
    +\underbrace{\sigma W(t_j)}_{\text{error}},
    \qquad j=0,\ldots,N.
  \]

  Since the Wiener process has mean zero, the expected log price is a line
  with intercept $\ln S_0$ and slope $\mug$. Least squares fits this line by
  minimizing the sum of squared differences from the observed log prices.

  \vspace{0.5em}
  The error $\sigma W(t_j)$ contains all shocks up to $t_j$, so errors at
  different times share shocks and are correlated. Their variance is
  $\sigma^2t_j$. These properties affect uncertainty in the fitted intercept
  and slope.
\end{frame}

\begin{frame}{Linear Regression: The Normal Equations}
  To fit the intercept and slope together, stack the $N+1$ log prices as
  $\hat{\mathbf X}\boldsymbol\theta+\boldsymbol\epsilon=\mathbf y$, where
  $\boldsymbol\theta=(\ln S_0,\mug)^{\top}$ contains the free intercept and slope:
  \[
    \hat{\mathbf X}=
    \begin{bmatrix}1&t_0\\1&t_1\\\vdots&\vdots\\1&t_N\end{bmatrix},
    \qquad
    \mathbf y=
    \begin{bmatrix}\ln S_{t_0}\\\ln S_{t_1}\\\vdots\\\ln S_{t_N}\end{bmatrix}.
  \]

  For full column rank (at least two distinct observation times), the
  least-squares solution is given by:
  \[
    \boxed{\hat{\boldsymbol\theta}
    =(\hat{\mathbf X}^{\top}\hat{\mathbf X})^{-1}\hat{\mathbf X}^{\top}\mathbf y.}
  \]
  The first component estimates the intercept. The second gives the mean-growth
  estimate $\hat{\mu}_g$, which we combine with volatility to
  recover $\hat\mu=\hat{\mu}_g+\hat\sigma^2/2$.
\end{frame}


\begin{frame}{Regression Intervals: Standard Errors}
  For comparison with Brownian errors, assume independent normal errors
  with mean zero and common variance $\sigma_\epsilon^2$.

  \vspace{0.4em}
  For $n=N+1$ observations and $p=2$ fitted parameters, the residuals and
  error variance estimate are:
  \[
    \mathbf r=\mathbf y-\hat{\mathbf X}\hat{\boldsymbol\theta},
    \qquad
    \hat\sigma_\epsilon^2=\frac{\|\mathbf r\|_2^2}{n-p}.
  \]
  With $n-p$ degrees of freedom after fitting, the standard error is given by:
  \[
    \boxed{\mathrm{SE}(\hat\theta_j)
      =\sqrt{\hat\sigma_\epsilon^2
        \left[(\hat{\mathbf X}^{\top}\hat{\mathbf X})^{-1}\right]_{jj}},
      \qquad j=1,2.}
  \]
  Standard error measures uncertainty in the fitted intercept or mean growth
  rate, in the parameter's units. Smaller values indicate greater precision.
\end{frame}

\begin{frame}{Regression Intervals: Confidence Intervals}
  We use the standard error to set an interval around each estimate.
  Under the independent normal-error assumptions, a confidence interval
  with level $1-\alpha$ for $\theta_j$ is given by:
  \[
    \boxed{\hat\theta_j\pm t_{1-\alpha/2,\nu}\,
      \mathrm{SE}(\hat\theta_j),\qquad \nu=n-p.}
  \]

  Here, $t_{1-\alpha/2,\nu}$ is the Student $t$ quantile with $\nu$ degrees of
  freedom. For a 95\% confidence interval, set $\alpha=0.05$, giving the
  multiplier $t_{0.975,\nu}$. Here, $j=1$ selects the intercept and
  $j=2$ mean growth.

  \vspace{0.6em}
  The \href{https://github.com/varnerlab/CHEME-5660-CourseRepository-Fall-2026/blob/main/lectures/week-4/L4b/advanced/drift-uncertainty/CHEME-5660-L4b-Advanced-DriftUncertainty-Fall-2026.ipynb}{advanced drift-uncertainty example}
  develops parameter uncertainty under GBM and examines its effect on target
  probabilities. Let's apply our estimates and ordinary regression
  intervals to historical share-price data.

  \slidenote{Example:}{\href{https://github.com/varnerlab/CHEME-5660-CourseRepository-Fall-2026/blob/main/lectures/week-4/L4b/CHEME-5660-L4b-Example-Parameters-SAGBM-Fall-2026.ipynb}{Estimating GBM Parameters}}
\end{frame}

\begin{frame}{Stylized Facts: Heavy-Tailed Growth Rates}
  Does fitted GBM reproduce L3a's patterns:
  heavy tails, little linear autocorrelation, and volatility clustering?
  Recall the model's growth-rate law:
  \[
    g_j=\mug+\frac{\sigma}{\sqrt{\Delta t}}Z_j,
    \qquad Z_j\overset{\mathrm{iid}}{\sim}\mathcal N(0,1).
  \]

  With constant parameters and independent normal draws, one-step growth
  rates are normal.

  \vspace{0.5em}
  \hilite{Heavy tails}: observed growth rates have heavier tails than a
  normal distribution. The model therefore assigns too little probability
  to extreme movements.

  \vspace{0.5em}
  The lognormal price has a long right tail, but growth rates depend on
  the log price ratio. A lognormal price does not give heavy-tailed growth rates.
\end{frame}

\begin{frame}{Stylized Facts: Little Linear Autocorrelation}
  From the growth-rate law, $g_j=\mug+(\sigma/\sqrt{\Delta t})Z_j$ with
  independent shocks. The autocovariance between growth rates $k\ge1$
  observation intervals apart is given by:
  \[
    \Cov(g_j,g_{j+k})
    =\frac{\sigma^2}{\Delta t}\,\Cov(Z_j,Z_{j+k})=0,
  \]
  because $Z_j$ and $Z_{j+k}$ are independent. With
  $\Var(g_j)=\sigma^2/\Delta t$, the autocorrelation function is:
  \[
    \boxed{\rho_g(k)=\frac{\Cov(g_j,g_{j+k})}{\Var(g_j)}=0,\qquad k\ge1.}
  \]

  This agrees with the weak linear dependence observed in signed growth
  rates. The model assumes independence, which is stronger than zero
  autocorrelation. Let's examine what independence implies for growth-rate magnitudes.
\end{frame}

\begin{frame}{Stylized Facts: Volatility Clustering}
  Volatility clustering concerns the \hilite{magnitudes} of growth rates.
  Large absolute movements tend to follow large movements, while small
  movements tend to follow small ones. This produces periods of high and
  low volatility in observed data.

  \vspace{0.7em}
  Under GBM, $\sigma$ is constant and shocks are independent. A large absolute
  growth rate therefore does not make another large movement more likely at
  the next step. The model cannot reproduce volatility clustering.

  \vspace{0.7em}
  Across these three patterns, GBM agrees with weak linear autocorrelation
  but misses heavy tails and volatility clustering. Let's
  use its analytical price distribution to evaluate a scheduled trade.
\end{frame}

\begin{frame}{GBM Trade Rule}
  Consider a scheduled trade: buy $n_0>0$ shares at $S_0$
  at time $0$ and sell at $S_T$ at time $T=N\Delta t$. Assume no dividends
  or trading costs.

  \vspace{0.4em}
  Let $g_y$ be the constant, continuously compounded benchmark growth rate
  (year$^{-1}$) corresponding to yield $y$. The NPV is given by:
  \[
    \operatorname{NPV}(g_y,T)
    =\underbrace{-n_0S_0}_{\text{purchase today}}
    +\underbrace{n_0S_Te^{-g_yT}}_{\text{discounted sale}}.
  \]

  Dividing by the initial investment gives the \hilite{scaled NPV}, our
  discounted fractional return:
  \[
    \rho_T=\frac{\operatorname{NPV}(g_y,T)}{n_0S_0}
    =\Bigl(\frac{S_T}{S_0}\Bigr)e^{-g_yT}-1.
  \]

  We now use GBM to describe the uncertain sale price $S_T$.
\end{frame}

\begin{frame}{Substitute the GBM Solution}
  Using mean growth rate $\mug=\mu-\sigma^2/2$, substitute the GBM price into
  the scaled NPV:
  \[
    \boxed{\rho_T=\exp\!\left[(\mug-g_y)T+\sigma\sqrt T\,Z\right]-1,
      \qquad Z\sim\mathcal N(0,1).}
  \]
  Discounting subtracts the benchmark growth rate $g_y$ from the mean growth
  rate $\mug$, while leaving the random shock unchanged. The quantity
  $1+\rho_T$ is lognormal, so:
  \[
    \ln(1+\rho_T)\sim\mathcal N\!\left((\mug-g_y)T,\sigma^2T\right).
  \]
  The scaled NPV is a lognormal variable shifted down by one, with
  $\rho_T>-1$.

  \vspace{0.5em}
  The benchmark affects the return through $g_yT$. When $|g_y|T\ll1$,
  discounting is small and $\rho_T\approx S_T/S_0-1$. We use the exact
  discounted expression to calculate the target probability.
\end{frame}

\begin{frame}{Terminal Target Probability}
  Choose a target scaled NPV $\rho_\star>-1$. Adding one and taking the
  logarithm preserves the strict target condition:
  \[
    \rho_T>\rho_\star
    \quad\Longleftrightarrow\quad
    \ln(1+\rho_T)>\ln(1+\rho_\star).
  \]
  Substituting the GBM expression gives:
  \[
    (\mug-g_y)T+\sigma\sqrt T\,Z>\ln(1+\rho_\star).
  \]
  Subtracting $(\mug-g_y)T$ and dividing by $\sigma\sqrt T>0$ gives a
  threshold for the standard normal random variable $Z$:
  \[
    Z>z_\star,\qquad
    \boxed{z_\star=\frac{\ln(1+\rho_\star)-(\mug-g_y)T}{\sigma\sqrt T}.}
  \]
  Thus $Z>z_\star$ means the trade exceeds its target at the scheduled sale.
  Selling when a boundary is first reached requires a separate calculation.
\end{frame}

\begin{frame}{The Closed-Form Target Probability}
  The standard normal cumulative distribution function $\Phi(z_\star)$ gives
  the probability that $Z\le z_\star$. Subtracting from one gives the
  probability that $Z>z_\star$:
  \[
    \boxed{\Pmeasure(\rho_T>\rho_\star)
      =1-\Phi(z_\star)
      =1-\Phi\!\left(
        \frac{\ln(1+\rho_\star)-(\mug-g_y)T}{\sigma\sqrt T}\right).}
  \]

  \vspace{0.4em}
  \begin{itemize}\setlength{\itemsep}{0.5em}
    \item For a chosen target, we evaluate $z_\star$ from estimated mean growth
      and volatility, the holding period, and benchmark growth rate.
    \item The formula applies for $T>0$, $\sigma>0$, and $\rho_\star>-1$.
    \item A target of $-1$ or below is exceeded with probability one because
      the scaled NPV always exceeds $-1$ under GBM.
  \end{itemize}
  \slidenote{Example:}{\href{https://github.com/varnerlab/CHEME-5660-CourseRepository-Fall-2026/blob/main/lectures/week-4/L4b/CHEME-5660-L4b-Example-GBM-NPV-TradeRule-Fall-2026.ipynb}{GBM Terminal-Target Probabilities}}
\end{frame}

% ==== optional advanced material =============================
\begin{frame}{Optional Advanced Material}
  \textbf{\ink{Advanced:}}
  \href{https://github.com/varnerlab/CHEME-5660-CourseRepository-Fall-2026/blob/main/lectures/week-4/L4b/advanced/lattice-limit/CHEME-5660-L4b-Advanced-LatticeToGBM-Fall-2026.ipynb}{From the Binomial Lattice to GBM}

  Calibrate the lattice to GBM and track convergence of its terminal distribution and target probability.

  \vspace{0.45em}
  \textbf{\ink{Advanced:}}
  \href{https://github.com/varnerlab/CHEME-5660-CourseRepository-Fall-2026/blob/main/lectures/week-4/L4b/advanced/first-passage/CHEME-5660-L4b-Advanced-FirstPassage-GBM-Fall-2026.ipynb}{First-Passage Rules Under GBM}

  Compare daily and continuous monitoring for single- and two-barrier exit rules.

  \vspace{0.45em}
  \textbf{\ink{Advanced:}}
  \href{https://github.com/varnerlab/CHEME-5660-CourseRepository-Fall-2026/blob/main/lectures/week-4/L4b/advanced/drift-uncertainty/CHEME-5660-L4b-Advanced-DriftUncertainty-Fall-2026.ipynb}{Drift Uncertainty}

  Examine how mean-growth precision depends on the observation period and affects a target probability.

  \vspace{0.45em}
  \textbf{\ink{Advanced:}}
  \href{https://github.com/varnerlab/CHEME-5660-CourseRepository-Fall-2026/blob/main/lectures/week-4/L4b/advanced/monte-carlo/CHEME-5660-L4b-Advanced-MonteCarlo-TargetProbability-Fall-2026.ipynb}{Monte Carlo Target Probabilities}

  Compare exact GBM and Euler--Maruyama simulation, quantify sampling uncertainty,
  and introduce variance reduction techniques.

  \vspace{0.4em}
  \dimmed{Optional; none of these notebooks is a prerequisite for L5a.}
\end{frame}

\begin{frame}{Summary}
  Today we developed a continuous-time share-price model, estimated its
  parameters, and used its price distribution to evaluate a stock trade.

  \vspace{0.3em}
  {\large\textbf{\ink{Key takeaways}}}
  \vspace{0.2em}
  \begin{itemize}\setlength{\itemsep}{0.35em}
    \item \hilite{Continuous-time price modeling}: We used the exact GBM solution
      to describe lognormal prices and simulate paths. It captures zero growth-rate
      autocorrelation but misses heavy tails and volatility clustering.
    \item \hilite{Mean growth, drift, and volatility}: We estimated mean growth
      from averages and log-price regression, and volatility from growth-rate
      dispersion. The half-variance correction recovers drift: $\mu=\mug+\sigma^2/2$.
    \item \hilite{Terminal target probabilities}: We used the GBM price
      distribution to calculate the probability of exceeding a target scaled
      NPV at the scheduled sale time, accounting for the holding period and
      benchmark growth rate.
  \end{itemize}

  \vspace{0.4em}
  Next, we extend GBM to correlated assets and introduce portfolio covariance.
\end{frame}

%\transitionframe{The GBM price distribution lets us calculate the probability that a scheduled trade exceeds a target discounted return.}

\end{document}
